\documentclass[a4paper,11pt,oneside,enabledeprecatedfontcommands]{scrbook}

% Pengaturan Font dan Bahasa untuk LuaLaTeX
\usepackage{fontspec}
\usepackage[indonesian]{babel}
\usepackage{amsmath,amssymb,amsfonts}
\usepackage{graphicx}
\usepackage[space]{grffile} % Mendukung nama file gambar yang mengandung spasi seperti "Tugas 4-001.png"
\usepackage{adjustbox}      % Mencegah rumus panjang keluar dari margin kanan
\usepackage{tikz}           % Pengganti ilustrasi C.a.R. jika file eksternal tidak tersedia
\usepackage{xcolor}
\usepackage{listings}
\usepackage{fancyvrb}
\usepackage{geometry}
\usepackage{enumitem}
\usepackage{hyperref}

\geometry{
    a4paper,
    top=2.5cm,
    bottom=2.5cm,
    left=2.5cm,
    right=2.5cm
}

\setlist[itemize]{label=-, topsep=2pt, partopsep=0pt, itemsep=2pt}
\setcounter{tocdepth}{1}

% Lokasi folder gambar
\graphicspath{{./}{images/}}

% Definisi warna ala Euler Math Toolbox (EMT)
\definecolor{eulerpromptcolor}{rgb}{0.6,0.0,0.0}
\definecolor{euleroutputcolor}{rgb}{0.15,0.15,0.45}
\definecolor{eulerbg}{rgb}{0.97,0.97,0.98}
\definecolor{eulerborder}{rgb}{0.82,0.85,0.90}

\hypersetup{
    colorlinks=true,
    linkcolor=euleroutputcolor,
    urlcolor=eulerpromptcolor,
    pdftitle={Menggambar Plot 3D dengan EMT}
}

% Pengaturan Listings untuk Kode EMT (rata kiri tanpa indentasi)
\lstset{
    basicstyle=\ttfamily\small,
    breaklines=true,
    breakatwhitespace=false,
    breakindent=0pt,
    breakautoindent=false,
    columns=fixed,
    basewidth=0.52em,
    keepspaces=true,
    showstringspaces=false,
    frame=leftline,
    framerule=1.5pt,
    framexleftmargin=6pt,
    xleftmargin=8pt,
    xrightmargin=0pt,
    resetmargins=true,
    rulecolor=\color{eulerborder},
    backgroundcolor=\color{eulerbg},
    aboveskip=6pt,
    belowskip=4pt,
    literate={°}{{\textdegree}}1
}

% Lingkungan khusus pengganti paket eumat agar kompatibel penuh dengan LuaLaTeX
\newenvironment{eulernotebook}{}{}

\newcommand{\eulerheading}[1]{\chapter{#1}}
\newcommand{\eulersubheading}[1]{\section{#1}}

\newenvironment{eulercomment}{\par\medskip\noindent}{\par\medskip}

% Otomatis menyesuaikan rumus yang terlalu lebar agar tidak terpotong margin kanan
\newenvironment{eulerformula}{%
    \par\medskip\noindent
    \renewcommand{\[}{\begin{center}\begin{adjustbox}{max width=\textwidth}$\displaystyle}%
    \renewcommand{\]}{$\end{adjustbox}\end{center}}%
}{\par\medskip}

\lstnewenvironment{eulerprompt}{
    \lstset{basicstyle=\ttfamily\small\color{eulerpromptcolor}}
}{}

\lstnewenvironment{euleroutput}{
    \lstset{
        basicstyle=\ttfamily\footnotesize\color{euleroutputcolor},
        backgroundcolor=\color{white},
        rulecolor=\color{lightgray}
    }
}{}

\lstnewenvironment{eulerudf}{
    \lstset{basicstyle=\ttfamily\small\color{teal!70!black}}
}{}

\lstnewenvironment{eulerttcomment}{
    \lstset{basicstyle=\ttfamily\small\color{darkgray}}
}{}

% Perintah penampil gambar yang menjaga proporsi rumus kecil (#1 < 10) maupun plot grafik (#1 >= 10)
\newcommand{\eulerimg}[2]{%
    \ifnum#1>0
        \begin{center}
            \IfFileExists{#2}{%
                \ifnum#1<10
                    \includegraphics[width=0.88\textwidth,height=#1\dimexpr 1.75ex\relax,keepaspectratio]{#2}%
                \else
                    \includegraphics[width=0.72\textwidth,height=#1\dimexpr 1.85ex\relax,keepaspectratio]{#2}%
                \fi
            }{%
                \fbox{\parbox{0.7\textwidth}{\centering\small\texttt{Gambar tidak ditemukan: \detokenize{#2}}\\ \footnotesize Pastikan folder \texttt{images/} berada satu direktori dengan file \texttt{.tex}.}}%
            }%
        \end{center}
    \fi
}

\title{\textbf{Menggambar Plot 3D dengan EMT}}
\subtitle{Panduan Visualisasi Permukaan, Kontur, Plot Implisit, dan Povray pada Euler Math Toolbox}
\author{Catatan Praktikum dan Latihan Plot 3D}
\date{\today}

\begin{document}

\frontmatter
\maketitle
\tableofcontents

\mainmatter
\begin{eulernotebook}

\eulerheading{Menggambar Plot 3D dengan EMT}
\begin{eulercomment}
Bagian ini merupakan pengantar pembuatan plot 3D di Euler untuk memvisualisasikan fungsi dua variabel.

Euler merender grafik fungsi menggunakan algoritma pengurutan untuk menyembunyikan elemen di latar belakang. Secara umum, Euler menerapkan proyeksi terpusat. Tampilan bawaan mengarah dari kuadran x-y positif
menuju titik asal (0,0,0). Pengaturan sudut 0° mengarahkan pandangan sejajar sumbu-y. Sudut dan ketinggian pandangan dapat disesuaikan kebutuhan.

Euler mampu memplot:
\begin{itemize}
    \item Permukaan dengan efek bayangan, garis kontur, atau rentang ketinggian.
    \item Kumpulan titik.
    \item Kurva parametrik.
    \item Permukaan implisit.
\end{itemize}

Plot 3D suatu fungsi dipanggil menggunakan perintah \texttt{plot3d}. Cara paling sederhana adalah memplot ekspresi matematika dengan variabel \texttt{x} dan \texttt{y}. Parameter \texttt{r} digunakan untuk mengatur rentang area plot di sekitar
titik asal (0,0).
\end{eulercomment}
\begin{eulerprompt}
>aspect(1.5); plot3d("x^2+sin(y)",-5,5,0,6*pi):
\end{eulerprompt}
\eulerimg{30}{images/Tugas 4-001.png}
\begin{eulerprompt}
>plot3d("x^2+x*sin(y)",-5,5,0,6*pi):
\end{eulerprompt}
\eulerimg{30}{images/Tugas 4-002.png}
\begin{eulercomment}
Modifikasi parameter grafik tersebut agar menghasilkan kurva lengkung atau melingkar pada bidang
horizontal maupun vertikal, lalu tentukan fungsi matematisnya.
\end{eulercomment}
\begin{eulerprompt}
>remvalue x, y
>plot3d("x*cos(y)", "x*sin(y)", "(x*cos(y))^2 + (x*cos(y))*sin(x*sin(y))", 0, 5, 0, 2*pi):
\end{eulerprompt}
\eulerimg{30}{images/Tugas 4-003.png}
\begin{eulerprompt}
>plot3d("sin(sqrt(x^2+y^2))*cos(y)", -5, 5, -5, 5):
\end{eulerprompt}
\eulerimg{30}{images/Tugas 4-004.png}
\begin{eulerprompt}
>plot3d("(4+cos(y))*cos(x)", "(4+cos(y))*sin(x)", "sin(y)+sin(x)", 0, 2*pi, 0, 2*pi):
\end{eulerprompt}
\eulerimg{30}{images/Tugas 4-005.png}

\eulerheading{Fungsi Dua Variabel}
\begin{eulercomment}
Pembuatan grafik fungsi dapat dilakukan menggunakan:
\begin{itemize}
    \item Ekspresi matematika dengan variabel \texttt{x} dan \texttt{y}.
    \item Pemanggilan nama fungsi dua variabel.
    \item Matriks data numerik.
\end{itemize}

Tampilan standar grafik berupa jaring kawat dengan warna berbeda di kedua sisinya. Nilai bawaan
interval jaring adalah 10, dan keseluruhan permukaan dibentuk oleh poligon 40x40. Pengaturan ini
dapat diubah:
\begin{itemize}
    \item \texttt{n=40}, \texttt{n=[40,40]}: Jumlah perulangan poligon pada setiap sumbu.
    \item \texttt{grid=10}, \texttt{grid=[10,10]}: Jumlah interval garis jaring pada setiap sumbu.
\end{itemize}

Contoh berikut menerapkan nilai parameter bawaan \texttt{n=40} dan \texttt{grid=10}.
\end{eulercomment}
\begin{eulerprompt}
>plot3d("x^2+y^2"):
\end{eulerprompt}
\eulerimg{30}{images/Tugas 4-006.png}
\begin{eulercomment}
Interaksi grafik dengan pengguna dilakukan dengan menambahkan parameter \texttt{>user}. Pengguna dapat
menggunakan tombol-tombol berikut:
\begin{itemize}
    \item Tombol panah (kiri, kanan, atas, bawah): Memutar sudut pandang.
    \item Tombol \texttt{+} dan \texttt{-}: Memperbesar atau memperkecil skala visual grafik.
    \item Tombol \texttt{a}: Menampilkan grafik dalam mode anaglif.
    \item Tombol \texttt{l}: Menyalakan atau mematikan rotasi sumber cahaya otomatis.
    \item Spasi: Mereset tampilan ke pengaturan sudut pandang awal.
    \item Return/Enter: Menutup sesi interaksi dan melanjutkan eksekusi program.
\end{itemize}
\end{eulercomment}
\begin{eulerprompt}
>plot3d("exp(-x^2+y^2)",>user, ...
>title="Putar dengan tombol panah (tekan return untuk selesai)"):
\end{eulerprompt}
\eulerimg{30}{images/Tugas 4-007.png}
\begin{eulercomment}
Batas koordinat grafik fungsi dapat diubah menggunakan parameter berikut:
\begin{itemize}
    \item \texttt{a,b}: Batas minimum dan maksimum nilai pada sumbu-x.
    \item \texttt{c,d}: Batas minimum dan maksimum nilai pada sumbu-y.
    \item \texttt{r}: Area simetris berupa bujur sangkar di sekitar titik asal (0,0).
    \item \texttt{n}: Jumlah titik persilangan atau interval dalam grid plot.
\end{itemize}

Parameter lain tersedia untuk memodifikasi skala fungsi atau rasio geometri grafik:
\begin{itemize}
    \item \texttt{fscale}: Memodifikasi rentang ukuran fungsi pada sumbu-z (defaultnya menggunakan \texttt{<fscale}).
    \item \texttt{scale}: Angka tunggal atau vektor 1x2 penentu rasio pelebaran arah x dan y.
    \item \texttt{frame}: Menentukan jenis bingkai pembatas grafik (nilai bawaan: 1).
\end{itemize}
\end{eulercomment}
\begin{eulerprompt}
>plot3d("exp(-(x^2+y^2)/5)",r=10,n=80,fscale=4,scale=1.2,frame=3,>user):
\end{eulerprompt}
\eulerimg{30}{images/Tugas 4-008.png}
\begin{eulercomment}
Proyeksi kamera diatur melalui parameter berikut:
\begin{itemize}
    \item \texttt{distance}: Jarak kamera virtual ke objek grafik.
    \item \texttt{zoom}: Faktor perbesaran tampilan.
    \item \texttt{angle}: Sudut rotasi sumbu-z dalam radian, diukur dari sumbu-y negatif.
    \item \texttt{height}: Sudut elevasi kamera dari bidang alas dalam radian.
\end{itemize}

Pengaturan kamera saat ini dapat diperiksa atau diubah dengan perintah \texttt{view()}. Perintah
ini mengembalikan matriks pengaturan kamera.
\end{eulercomment}
\begin{eulerprompt}
>view
\end{eulerprompt}
\begin{euleroutput}
[5,  2.6,  2,  0.4]
\end{euleroutput}
\begin{eulercomment}
Mengatur jarak kamera yang kecil menghilangkan perlunya pembesaran (zoom) tinggi,
menghasilkan efek proyeksi sudut lebar.

Contoh berikut menggunakan sudut (\texttt{angle}) pi/2 dan elevasi (\texttt{height}) 0, sehingga kamera
melihat tepat dari sumbu-y negatif. Label indikator untuk sumbu-y secara otomatis tidak
ditampilkan.
\end{eulercomment}
\begin{eulerprompt}
>plot3d("x^2+y",distance=3,zoom=1,angle=pi/2,height=0):
\end{eulerprompt}
\eulerimg{30}{images/Tugas 4-009.png}
\begin{eulercomment}
Secara standar, pusat pandangan mengarah ke titik tengah kubus pembatas. Pusat pandangan
ini dapat digeser dengan parameter \texttt{center}.
\end{eulercomment}
\begin{eulerprompt}
>plot3d("x^4+y^2",a=0,b=1,c=-1,d=1,angle=-20°,height=20°, ...
>center=[0.4,0,0],zoom=5):
\end{eulerprompt}
\eulerimg{30}{images/Tugas 4-010.png}
\begin{eulercomment}
Grafik 3D disesuaikan otomatis agar muat ke dalam kubus standar. Hal ini memudahkan
penyesuaian tanpa perlu mengatur jarak kamera atau zoom secara manual saat rentang skala
bervariasi. Penanda label angka akan mengindikasikan nilai absolut riil pada grafik.

Jika normalisasi rasio skala dimatikan (\texttt{<scale} atau \texttt{scale=false}), Anda perlu mengatur
nilai \texttt{distance} dan \texttt{center} secara manual agar geometri grafik tidak terpotong.
\end{eulercomment}
\begin{eulerprompt}
>plot3d("5*exp(-x^2-y^2)",r=2,<fscale,<scale,distance=13,height=50°, ...
>center=[0,0,-2],frame=3):
\end{eulerprompt}
\eulerimg{30}{images/Tugas 4-011.png}
\begin{eulercomment}
Proyeksi polar dapat diaktifkan dengan menambahkan parameter \texttt{>polar}. Fungsi matematika
harus mendefinisikan \texttt{r} terhadap \texttt{x} dan sudut terhadap \texttt{y}. Pengubahan skala fungsi polar
ditangani oleh \texttt{fscale}, sementara bidang horizontal disesuaikan dengan skala otomatis
bawaan.
\end{eulercomment}
\begin{eulerprompt}
>plot3d("1/(x^2+y^2+1)",r=5,>polar, ...
>fscale=2,>hue,n=100,zoom=4,>contour,color=blue):
\end{eulerprompt}
\eulerimg{30}{images/Tugas 4-012.png}
\begin{eulerprompt}
>function f(r) := exp(-r/2)*cos(r); ...
>plot3d("f(x^2+y^2)",>polar,scale=[1,1,0.4],r=pi,frame=3,zoom=4):
\end{eulerprompt}
\eulerimg{30}{images/Tugas 4-013.png}
\begin{eulercomment}
Parameter \texttt{rotate} memungkinkan pemutaran kurva 2D mengelilingi sumbu tertentu untuk
membentuk permukaan 3D:
\begin{itemize}
    \item \texttt{rotate=1} (atau \texttt{rotate=true}): Memutar kurva mengelilingi sumbu-x.
    \item \texttt{rotate=2}: Memutar kurva mengelilingi sumbu-z.
\end{itemize}
\end{eulercomment}
\begin{eulerprompt}
>plot3d("x^2+1",a=-1,b=1,rotate=true,grid=5):
\end{eulerprompt}
\eulerimg{30}{images/Tugas 4-014.png}
\begin{eulerprompt}
>plot3d("x^2+1",a=-1,b=1,rotate=2,grid=5):
\end{eulerprompt}
\eulerimg{30}{images/Tugas 4-015.png}
\begin{eulerprompt}
>plot3d("sqrt(25-x^2)",a=0,b=5,rotate=1):
\end{eulerprompt}
\eulerimg{30}{images/Tugas 4-016.png}
\begin{eulerprompt}
>plot3d("x*sin(x)",a=0,b=6pi,rotate=2):
\end{eulerprompt}
\eulerimg{30}{images/Tugas 4-017.png}
\begin{eulercomment}
Berikut adalah contoh menggabungkan tiga ekspresi matematika sekaligus dalam satu plot
3D.
\end{eulercomment}
\begin{eulerprompt}
>plot3d("x","x^2+y^2","y",r=2,zoom=3.5,frame=3):
\end{eulerprompt}
\eulerimg{30}{images/Tugas 4-018.png}

\eulerheading{Plot Kontur}
\begin{eulercomment}
Standarnya, Euler menampilkan jaring (grid) persegi pada plot. Grid ini dapat dihilangkan
untuk menonjolkan gradasi warna berdasarkan ketinggian (spectral hue). Euler juga
mendukung efek pencahayaan (shading) untuk mempertegas bentuk permukaan serta pembuatan
gambar anaglif tiga dimensi (merah/sian).
\begin{itemize}
    \item \texttt{>hue}: Mengaktifkan pewarnaan bayangan (shading) dan menyembunyikan garis jaring kawat.
    \item \texttt{>contour}: Menampilkan garis kontur ketinggian pada permukaan grafik.
    \item \texttt{level=...} (atau \texttt{levels}): Menentukan rentang atau nilai level untuk garis kontur.
\end{itemize}

Euler menggunakan pengaturan level otomatis (\texttt{level="auto"}) untuk
mendistribusikan garis kontur secara proporsional.

Nilai parameter dapat disesuaikan. Contoh berikut menggunakan grid 100x100, menyesuaikan
skala sumbu, serta mengubah sudut pandang kamera.
\end{eulercomment}
\begin{eulerprompt}
>plot3d("exp(-x^2-y^2)",r=2,n=100,level="thin", ...
>>contour,>spectral,fscale=1,scale=1.1,angle=45°,height=20°):
\end{eulerprompt}
\eulerimg{30}{images/Tugas 4-019.png}
\begin{eulerprompt}
>plot3d("exp(x*y)",angle=100°,>contour,color=green):
\end{eulerprompt}
\eulerimg{30}{images/Tugas 4-020.png}
\begin{eulercomment}
Pengaturan awal memberikan warna abu-abu polos pada shading. Skema warna dapat diubah
menggunakan palet bawaan:
\begin{itemize}
    \item \texttt{>spectral}: Menggunakan skema warna spektral standar.
    \item \texttt{color=...}: Memilih warna kustom atau skema warna khusus.
\end{itemize}

Contoh berikut meningkatkan jumlah titik kalkulasi (\texttt{n=100}) agar garis kontur dirender
lebih halus dan presisi.
\end{eulercomment}
\begin{eulerprompt}
>plot3d("x^2+y^2",>spectral,>contour,n=100):
\end{eulerprompt}
\eulerimg{30}{images/Tugas 4-021.png}
\begin{eulercomment}
Rentang dan nilai garis kontur pada permukaan grafik dapat ditentukan secara eksplisit.
Pada contoh berikut, tingkat kontur ditentukan dari -1 hingga 1 dengan interval 0.1.
\end{eulercomment}
\begin{eulerprompt}
>plot3d("x^2-y^2",0,5,0,5,level=-1:0.1:1,color=redgreen):
\end{eulerprompt}
\eulerimg{30}{images/Tugas 4-022.png}
\begin{eulercomment}
Contoh berikut menampilkan pita kontur yang lebih tebal pada interval nilai [-0.1, 0] dan
[0.9, 1]. Format nilai disajikan dalam matriks 2x2.

Parameter \texttt{grid=10} digunakan untuk membagi area menjadi 10 interval jaring.
\end{eulercomment}
\begin{eulerprompt}
>plot3d("x^2+y^3",level=[-0.1,0.9;0,1], ...
>>spectral,angle=30°,grid=10,contourcolor=gray):
\end{eulerprompt}
\eulerimg{30}{images/Tugas 4-023.png}
\begin{eulercomment}
Contoh berikut memplot kurva tempat titik-titik yang memenuhi persamaan implisit:
\end{eulercomment}
\begin{eulerformula}
\[
f(x,y) = x^y-y^x = 0
\]
\end{eulerformula}
\begin{eulercomment}
Plot ini ditampilkan khusus untuk level tunggal pada ketinggian zero set (\texttt{level=0}).
\end{eulercomment}
\begin{eulerprompt}
>plot3d("x^y-y^x",level=0,a=0,b=6,c=0,d=6,contourcolor=red,n=100):
\end{eulerprompt}
\eulerimg{30}{images/Tugas 4-025.png}
\begin{eulercomment}
Proyeksi garis kontur juga dapat diproyeksikan ke alas bawah permukaan grafik menggunakan
parameter \texttt{>cp}. Warna kontur proyeksi dan jaraknya dari plot dapat dikustomisasi.
\end{eulercomment}
\begin{eulerprompt}
>plot3d("x^2+y^4",>cp,cpcolor=green,cpdelta=0.2):
\end{eulerprompt}
\eulerimg{30}{images/Tugas 4-026.png}
\begin{eulercomment}
Berikut adalah contoh berbagai gaya tampilan dan pewarnaan pada Euler. Bingkai objek
dapat dimatikan dengan atribut \texttt{<frame}, disertai penyesuaian skema warna serta
kepadatan jaring (grid).
\end{eulercomment}
\begin{eulerprompt}
>figure(2,2); ...
>expr="y^3-x^2"; ...
>figure(1); ...
>plot3d(expr,<frame,>cp,cpcolor=spectral); ...
>figure(2); ...
>plot3d(expr,<frame,>spectral,grid=10,cp=2); ...
>figure(3); ...
>plot3d(expr,<frame,>contour,color=gray,nc=5,cp=3,cpcolor=greenred); ...
>figure(4); ...
>plot3d(expr,<frame,>hue,grid=10,>transparent,>cp,cpcolor=gray); ...
>figure(0):
\end{eulerprompt}
\eulerimg{30}{images/Tugas 4-027.png}
\begin{eulercomment}
Skema warna spektral standar menggunakan indeks bernomor 1 hingga 9. Indeks ini dapat
dipanggil menggunakan atribut \texttt{spectral=i} atau \texttt{color=value}. Pilihan warna bawaan meliputi:
\begin{itemize}
    \item \texttt{spectral}: Gradasi warna dari biru ke merah.
    \item \texttt{white}: Warna dasar putih dengan shading.
    \item Kombinasi dua warna: \texttt{yellowblue}, \texttt{purplegreen}, \texttt{blueyellow}, \texttt{greenred}, \texttt{greenpurple}, \texttt{redgreen}.
\end{itemize}
\end{eulercomment}
\begin{eulerprompt}
>figure(3,3); ...
>for i=1:9; ...
>figure(i); plot3d("x^2+y^2",spectral=i,>contour,>cp,<frame,zoom=4); ...
>end; ...
>figure(0):
\end{eulerprompt}
\eulerimg{30}{images/Tugas 4-028.png}
\begin{eulercomment}
Arah sumber pencahayaan dapat diubah interaktif saat grafik ditampilkan dengan menekan
tombol \texttt{'l'} dan menggunakan tombol panah. Posisi arah pencahayaan juga dapat
ditentukan secara permanen melalui kode skrip:
\begin{itemize}
    \item \texttt{light}: Vektor arah pencahayaan virtual.
    \item \texttt{amb}: Intensitas cahaya ambien (berkisar dari 0 hingga 1).
\end{itemize}

Modul perenderan internal Euler tidak menghitung proyeksi bayangan jatuh (shadows). Untuk
merender bayangan objek yang fotorealistik, pengguna dapat memanfaatkan integrasi dengan
Povray.
\end{eulercomment}
\begin{eulerprompt}
>plot3d("-x^2-y^2", ...
>hue=true,light=[0,1,1],amb=0,user=true, ...
>title="Tekan 'l' dan tombol kursor (tekan return untuk keluar)"):
\end{eulerprompt}
\eulerimg{30}{images/Tugas 4-029.png}
\begin{eulercomment}
Atribut \texttt{color} menetapkan warna tunggal pada permukaan grafik. Warna garis indikator level
kontur juga dapat disesuaikan (\texttt{contourcolor}).
\end{eulercomment}
\begin{eulerprompt}
>plot3d("-x^2-y^2",color=rgb(0.2,0.2,0),hue=true,frame=false, ...
>zoom=3,contourcolor=red,level=-2:0.1:1,dl=0.01):
\end{eulerprompt}
\eulerimg{30}{images/Tugas 4-030.png}
\begin{eulercomment}
Menyetel \texttt{color=0} mengaktifkan efek warna pelangi (rainbow effect) pada shading permukaan
grafik.
\end{eulercomment}
\begin{eulerprompt}
>plot3d("x^2/(x^2+y^2+1)",color=0,hue=true,grid=10):
\end{eulerprompt}
\eulerimg{30}{images/Tugas 4-031.png}
\begin{eulercomment}
Permukaan poligon grafik dapat dibuat transparan agar struktur bagian dalam terlihat.
Fitur ini diaktifkan dengan atribut \texttt{>transparent}.
\end{eulercomment}
\begin{eulerprompt}
>plot3d("x^2+y^2",>transparent,grid=10,wirecolor=red):
\end{eulerprompt}
\eulerimg{30}{images/Tugas 4-032.png}

\eulerheading{Plot Implisit}
\begin{eulercomment}
Euler mendukung perenderan permukaan dari persamaan fungsi implisit $f(x,y,z)=0$. Algoritma
\texttt{plot3d} memotong bidang ruang untuk menampilkan permukaan zero set fungsi tersebut.
Titik-titik penyelesaian dari persamaan:
\end{eulercomment}
\begin{eulerformula}
\[
f(x,y,z) = 0
\]
\end{eulerformula}
\begin{eulercomment}
diproyeksikan pada irisan bidang x-y, x-z, atau y-z:
\begin{itemize}
    \item \texttt{implicit=1}: Memotong irisan paralel bidang y-z.
    \item \texttt{implicit=2}: Memotong irisan paralel bidang x-z.
    \item \texttt{implicit=4}: Memotong irisan paralel bidang x-y.
\end{itemize}

Kombinasi nilai dapat digunakan untuk menampilkan beberapa irisan sekaligus. Contoh
berikut memplot fungsi implisit:
\end{eulercomment}
\begin{eulerformula}
\[
M = \{ (x,y,z) : x^2+y^3+zy=1 \}
\]
\end{eulerformula}
\begin{eulerprompt}
>plot3d("x^2+y^3+z*y-1",r=5,implicit=3):
\end{eulerprompt}
\eulerimg{30}{images/Tugas 4-035.png}
\begin{eulerprompt}
>c=1; d=1;
>plot3d("((x^2+y^2-c^2)^2+(z^2-1)^2)*((y^2+z^2-c^2)^2+(x^2-1)^2)*((z^2+x^2-c^2)^2+(y^2-1)^2)-d",r=2,<frame,>implicit,>user):
\end{eulerprompt}
\eulerimg{30}{images/Tugas 4-036.png}
\begin{eulerprompt}
>plot3d("x^2+y^2+4*x*z+z^3",>implicit,r=2,zoom=2.5):
\end{eulerprompt}
\eulerimg{30}{images/Tugas 4-037.png}

\eulerheading{Memplot Data 3D}
\begin{eulercomment}
Selain persamaan fungsi, \texttt{plot3d} dapat menerima data numerik berupa tiga matriks koordinat
(x,y,z) untuk menampilkan objek 3D parametrik, serupa dengan fungsi \texttt{plot2d}. Anda dapat
memasukkan matriks x, y, dan z secara langsung, atau menggunakan fungsi parametrik
$fx(t,s)$, $fy(t,s)$, dan $fz(t,s)$.
\end{eulercomment}
\begin{eulerformula}
\[
\gamma(t,s) = (x(t,s),y(t,s),z(t,s))
\]
\end{eulerformula}
\begin{eulercomment}
Parameter $t$ dan $s$ dievaluasi dalam grid persegi untuk menghasilkan matriks koordinat x,
y, dan z yang membentuk permukaan 3D.

Penggunaan kalkulasi berarah matriks di Euler sangat direkomendasikan karena lebih
efisien dalam memproses data koordinat.

Contoh berikut memplot permukaan bola menggunakan parameter \texttt{t} (vektor baris) dan \texttt{s}
(vektor kolom). Wilayah kutub juga dapat diwarnai khusus menggunakan parameter \texttt{level} dan
\texttt{contourcolor}.
\end{eulercomment}
\begin{eulerprompt}
>t=linspace(0,2pi,180); s=linspace(-pi/2,pi/2,90)'; ...
>x=cos(s)*cos(t); y=cos(s)*sin(t); z=sin(s); ...
>plot3d(x,y,z,>hue, ...
>color=blue,<frame,grid=[10,20], ...
>values=s,contourcolor=red,level=[90°-24°;90°-22°], ...
>scale=1.4,height=50°):
\end{eulerprompt}
\eulerimg{30}{images/Tugas 4-039.png}
\begin{eulercomment}
Contoh berikut menampilkan grafik fungsi parametrik $z = f(t,s)$:
\end{eulercomment}
\begin{eulerprompt}
>t=-1:0.1:1; s=(-1:0.1:1)'; plot3d(t,s,t*s,grid=10):
\end{eulerprompt}
\eulerimg{30}{images/Tugas 4-040.png}
\begin{eulercomment}
Permukaan juga dapat didefinisikan secara implisit sebagai fungsi $x$ terhadap $y$ dan $z$:
\end{eulercomment}
\begin{eulerformula}
\[
x = y \, z
\]
\end{eulerformula}
\begin{eulerprompt}
>plot3d(t*s,t,s,angle=180°,grid=10):
\end{eulerprompt}
\eulerimg{30}{images/Tugas 4-042.png}
\begin{eulercomment}
Berbagai bentuk permukaan dapat dibuat dengan mendefinisikan persamaan parametriknya
secara tepat.

Tampilan berikut menunjukkan bola terdistorsi dengan shading. Permukaan bola standar
didefinisikan dengan persamaan:
\end{eulercomment}
\begin{eulerformula}
\[
\gamma(t,s) = (\cos(t)\cos(s),\sin(t)\sin(s),\cos(s))
\]
\end{eulerformula}
\begin{eulercomment}
dengan rentang parameter:
\end{eulercomment}
\begin{eulerformula}
\[
0 \le t \le 2\pi, \quad \frac{-\pi}{2} \le s \le \frac{\pi}{2}.
\]
\end{eulerformula}
\begin{eulercomment}
Bentuk bola ini kemudian didistorsi menggunakan faktor pengali $d(t,s)$:
\end{eulercomment}
\begin{eulerformula}
\[
d(t,s) = \frac{\cos(4t)+\cos(8s)}{4}.
\]
\end{eulerformula}
\begin{eulerprompt}
>t=linspace(0,2pi,320); s=linspace(-pi/2,pi/2,160)'; ...
>d=1+0.2*(cos(4*t)+cos(8*s)); ...
>plot3d(cos(t)*cos(s)*d,sin(t)*cos(s)*d,sin(s)*d,hue=1, ...
>light=[1,0,1],frame=0,zoom=5):
\end{eulerprompt}
\eulerimg{30}{images/Tugas 4-046.png}
\begin{eulercomment}
Euler juga dapat memplot kumpulan titik 3D (point cloud) dari tiga vektor data koordinat.

Opsi gaya tampilan yang tersedia pada \texttt{plot2d} dapat digunakan pada \texttt{plot3d} dengan
menambahkan parameter \texttt{points=true}.
\end{eulercomment}
\begin{eulerprompt}
>n=500; ...
>plot3d(normal(1,n),normal(1,n),normal(1,n),points=true,style="."):
\end{eulerprompt}
\eulerimg{30}{images/Tugas 4-047.png}
\begin{eulercomment}
Kurva 3D dapat dibuat dari urutan titik koordinat dengan mengaktifkan parameter \texttt{>wire}.
Data titik koordinat dihitung terlebih dahulu sebelum diplot.
\end{eulercomment}
\begin{eulerprompt}
>t=linspace(0,8pi,500); ...
>plot3d(sin(t),cos(t),t/10,>wire,zoom=3):
\end{eulerprompt}
\eulerimg{30}{images/Tugas 4-048.png}
\begin{eulerprompt}
>t=linspace(0,4pi,1000); plot3d(cos(t),sin(t),t/2pi,>wire, ...
>linewidth=3,wirecolor=blue):
\end{eulerprompt}
\eulerimg{30}{images/Tugas 4-049.png}
\begin{eulerprompt}
>X=cumsum(normal(3,100)); ...
>plot3d(X[1],X[2],X[3],>anaglyph,>wire):
\end{eulerprompt}
\eulerimg{30}{images/Tugas 4-050.png}
\begin{eulercomment}
Euler dapat menampilkan grafik dalam mode anaglif 3D yang dapat dilihat menggunakan
kacamata filter merah/sian.
\end{eulercomment}
\begin{eulerprompt}
>plot3d("x^2+y^3",>anaglyph,>contour,angle=30°):
\end{eulerprompt}
\eulerimg{30}{images/Tugas 4-051.png}
\begin{eulercomment}
Perbedaan ketinggian permukaan grafik dapat ditonjolkan dengan pewarnaan spektral
bergradasi.
\end{eulercomment}
\begin{eulerprompt}
>plot3d("x^2*y^3-y",>spectral,>contour,zoom=3.2):
\end{eulerprompt}
\eulerimg{30}{images/Tugas 4-052.png}
\begin{eulercomment}
Permukaan terparameterisasi dirender dengan memasukkan matriks nilai koordinat x, y, dan
z yang dihasilkan dari jaring nilai parameter.

Contoh berikut memplot permukaan pita Mobius menggunakan parameter \texttt{u} dan \texttt{v}.
\end{eulercomment}
\begin{eulerprompt}
>u=linspace(-1,1,10); v=linspace(0,2*pi,50)'; ...
>X=(3+u*cos(v/2))*cos(v); Y=(3+u*cos(v/2))*sin(v); Z=u*sin(v/2); ...
>plot3d(X,Y,Z,>anaglyph,<frame,>wire,scale=2.3):
\end{eulerprompt}
\eulerimg{30}{images/Tugas 4-053.png}
\begin{eulercomment}
Efek kedalaman visual 3D pada grafik berikut akan terlihat optimal jika dilihat
menggunakan kacamata anaglif merah/sian.
\end{eulercomment}
\begin{eulerprompt}
>u:=linspace(-pi,pi,160); v:=linspace(-pi,pi,400)'; ...
>x:=(4*(1+.25*sin(3*v))+cos(u))*cos(2*v); ...
>y:=(4*(1+.25*sin(3*v))+cos(u))*sin(2*v); ...
>z=sin(u)+2*cos(3*v); ...
>plot3d(x,y,z,frame=0,scale=1.5,hue=1,light=[1,0,-1],zoom=2.8,>anaglyph):
\end{eulerprompt}
\eulerimg{30}{images/Tugas 4-054.png}

\eulerheading{Plot Statistik}
\begin{eulercomment}
Diagram batang 3D (bar plot) dapat dibuat menggunakan \texttt{plot3d} dengan memasukkan data
matriks:
\begin{itemize}
    \item \texttt{x}: Vektor baris bernilai $n+1$ sebagai batas sel sumbu-x.
    \item \texttt{y}: Vektor kolom bernilai $n+1$ sebagai batas sel sumbu-y.
    \item \texttt{z}: Matriks ukuran $n \times n$ yang berisi nilai tinggi batang.
\end{itemize}

Ukuran matriks \texttt{z} dapat lebih besar dari batas \texttt{x} dan \texttt{y}, tetapi hanya bagian $n \times n$ yang akan
diplot.

Pada contoh berikut, nilai \texttt{x} dan \texttt{y} disesuaikan terlebih dahulu agar posisi batang berada
tepat di tengah koordinat yang sesuai.
\end{eulercomment}
\begin{eulerprompt}
>x=-1:0.1:1; y=x'; z=x^2+y^2; ...
>xa=(x|1.1)-0.05; ya=(y_1.1)-0.05; ...
>plot3d(xa,ya,z,bar=true):
\end{eulerprompt}
\eulerimg{30}{images/Tugas 4-055.png}
\begin{eulercomment}
Permukaan grafik dapat dipisahkan menjadi beberapa segmen menggunakan parameter
\texttt{disconnect}.
\end{eulercomment}
\begin{eulerprompt}
>x=-1:0.1:1; y=x'; z=x+y; d=zeros(size(x)); ...
>plot3d(x,y,z,disconnect=2:2:20):
\end{eulerprompt}
\eulerimg{30}{images/Tugas 4-056.png}
\begin{eulercomment}
Saat memplot matriks data \texttt{M} dari berkas, skala nilai sumbu-z akan dinormalisasi ke
interval $[-1,1]$. Pengaturan skala otomatis pada sumbu-z ini dapat diaktifkan dengan
menambahkan parameter \texttt{>zscale}, atau disesuaikan dengan skala kustom (\texttt{scale}).
\end{eulercomment}
\begin{eulerprompt}
>i=1:20; j=i'; ...
>plot3d(i*j^2+100*normal(20,20),>zscale,scale=[1,1,1.5],angle=-40°,zoom=1.8):
\end{eulerprompt}
\eulerimg{30}{images/Tugas 4-057.png}
\begin{eulerprompt}
>Z=intrandom(5,100,6); v=zeros(5,6); ...
>loop 1 to 5; v[#]=getmultiplicities(1:6,Z[#]); end; ...
>columnsplot3d(v',scols=1:5,ccols=[1:5]):
\end{eulerprompt}
\eulerimg{30}{images/Tugas 4-058.png}

\eulerheading{Permukaan Benda Putar}
\begin{eulerprompt}
>plot2d("(x^2+y^2-1)^3-x^2*y^3",r=1.3, ...
>style="#",color=red,<outline, ...
>level=[-2;0],n=100):
\end{eulerprompt}
\eulerimg{30}{images/Tugas 4-059.png}
\begin{eulerprompt}
>ekspresi &= (x^2+y^2-1)^3-x^2*y^3; $ekspresi
\end{eulerprompt}
\begin{eulerformula}
\[
\left(y^2+x^2-1\right)^3-x^2\,y^3
\]
\end{eulerformula}
\begin{eulercomment}
Kurva bentuk hati (heart curve) dapat diputar mengelilingi sumbu-y untuk membentuk benda
putar 3D. Persamaan kurvanya adalah:
\end{eulercomment}
\begin{eulerformula}
\[
f(x,y)=(x^2+y^2-1)^3-x^2.y^3.
\]
\end{eulerformula}
\begin{eulercomment}
Untuk memperoleh bentuk parametriknya, dilakukan substitusi koordinat polar:
\end{eulercomment}
\begin{eulerformula}
\[
x=r.cos(a),\quad y=r.sin(a).
\]
\end{eulerformula}
\begin{eulerprompt}
>function fr(r,a) &= ekspresi with [x=r*cos(a),y=r*sin(a)] | trigreduce; $fr(r,a)
\end{eulerprompt}
\begin{eulerformula}
\[
\left(r^2-1\right)^3+\frac{\left(\sin \left(5\,a\right)-\sin \left(  3\,a\right)-2\,\sin a\right)\,r^5}{16}
\]
\end{eulerformula}
\begin{eulercomment}
Nilai jari-jari $r$ dicari menggunakan fungsi \texttt{bisect} berdasarkan sudut $a$. Hasil kurva 3D
parametrik kemudian dipromosikan dan diplot sebagai permukaan benda putar.
\end{eulercomment}
\begin{eulerprompt}
>function map f(a) := bisect("fr",0,2;a); ...
>t=linspace(-pi/2,pi/2,100); r=f(t); ...
>s=linspace(pi,2pi,100)'; ...
>plot3d(r*cos(t)*sin(s),r*cos(t)*cos(s),r*sin(t), ...
>>hue,<frame,color=red,zoom=4,amb=0,max=0.7,grid=12,height=50°):
\end{eulerprompt}
\eulerimg{30}{images/Tugas 4-064.png}
\begin{eulercomment}
Benda putar juga dapat dibuat dengan memutar grafik implisit mengelilingi sumbu-z.
Persamaan fungsi didefinisikan sebagai berikut:
\end{eulercomment}
\begin{eulerprompt}
>function f(x,y,z) ...
\end{eulerprompt}
\begin{eulerudf}
r=x^2+y^2;
return (r+z^2-1)^3-r*z^3;
endfunction
\end{eulerudf}
\begin{eulerprompt}
>plot3d("f(x,y,z)", ...
>xmin=0,xmax=1.2,ymin=-1.2,ymax=1.2,zmin=-1.2,zmax=1.4, ...
>implicit=1,angle=-30°,zoom=2.5,n=[10,100,60],>anaglyph):
\end{eulerprompt}
\eulerimg{30}{images/Tugas 4-065.png}

\eulerheading{Plot 3D Khusus}
\begin{eulercomment}
Fungsi \texttt{plot3d} memudahkan pembuatan grafik 3D secara cepat, tetapi tidak semua kebutuhan
khusus dapat dipenuhi. Pengguna dapat membuat prosedur kustom untuk memplot objek khusus
sesuai kebutuhan.

EMT bukan merupakan program ray tracer 3D khusus, tetapi dapat menggabungkan beberapa
objek plot 3D sekaligus. Contoh berikut memvisualisasikan permukaan paraboloida beserta
bidang singgungnya (tangent plane).
\end{eulercomment}
\begin{eulerprompt}
>function myplot ...
\end{eulerprompt}
\begin{eulerudf}
y=-1:0.01:1; x=(-1:0.01:1)';
plot3d(x,y,0.2*(x-0.1)/2,<scale,<frame,>hue, ..
hues=0.5,>contour,color=orange);
h=holding(1);
plot3d(x,y,(x^2+y^2)/2,<scale,<frame,>contour,>hue);
holding(h);
endfunction
\end{eulerudf}
\begin{eulercomment}
Prosedur \texttt{framedplot()} digunakan untuk mengatur bingkai kanvas serta menentukan batas
koordinat dan sudut pandang kamera.
\end{eulercomment}
\begin{eulerprompt}
>framedplot("myplot",[-1,1,-1,1,0,1],height=0,angle=-30°, ...
>center=[0,0,-0.7],zoom=3):
\end{eulerprompt}
\eulerimg{30}{images/Tugas 4-066.png}
\begin{eulercomment}
Proyeksi garis kontur pada alas bidang (contour plane) dihitung terpisah. Fungsi \texttt{plot3d()}
bekerja pada tampilan jendela penuh (\texttt{fullwindow()}), sedangkan \texttt{plotcontourplane()}
menggambar garis kontur sesuai batas bingkai yang ditentukan.
\end{eulercomment}
\begin{eulerprompt}
>x=-1:0.02:1.1; y=x'; z=x^2-y^4;
>function myplot (x,y,z) ...
\end{eulerprompt}
\begin{eulerudf}
zoom(2);
wi=fullwindow();
plotcontourplane(x,y,z,level="auto",<scale);
plot3d(x,y,z,>hue,<scale,>add,color=white,level="thin");
window(wi);
reset();
endfunction
\end{eulerudf}
\begin{eulerprompt}
>myplot(x,y,z):
\end{eulerprompt}
\eulerimg{25}{images/Tugas 4-067.png}

\eulerheading{Animasi}
\begin{eulercomment}
EMT dapat menghasilkan animasi 3D dengan merekam urutan gambar (frame) dari perubahan
parameter grafik secara berkala.

Fungsi \texttt{rotate()} digunakan untuk memutar sudut pandang grafik, meredraw tampilan, dan
menyimpan setiap bingkai menggunakan \texttt{addpage()} hingga menghasilkan animasi pemutaran
objek.

Pelajari kode sumber fungsi \texttt{rotate()} untuk memahami alur pembuatan animasi secara rinci.
\end{eulercomment}
\begin{eulerprompt}
>function testplot () := plot3d("x^2+y^3"); ...
>rotate("testplot"); testplot():
\end{eulerprompt}
\eulerimg{25}{images/Tugas 4-068.png}

\eulerheading{Menggambar dengan Povray}
\begin{eulercomment}
Pustaka \texttt{povray.e} memungkinkan pengguna mengekspor data EMT ke format berkas Povray. Hasil
pencitraan ray tracing dari Povray jauh lebih halus dan fotorealistik dibandingkan
tampilan grafik standar EMT.

Pastikan program Povray sudah terpasang di sistem komputer Anda (\url{http://www.povray.org/})
dan jalurnya terdaftar dalam environment path. Tentukan lokasi berkas eksekusi
\texttt{pvengine.exe} pada variabel \texttt{defaultpovray}.

EMT akan menyimpan deskripsi objek ke berkas \texttt{current.pov} di direktori \texttt{eulerhome()}.
Selanjutnya, EMT memanggil Povray untuk merender berkas tersebut menjadi berkas gambar
PNG yang akan ditampilkan di dokumen. Berkas sementara POV dapat dibersihkan menggunakan
perintah \texttt{povclear()}.

Fungsi \texttt{pov3d} memiliki sintaks yang serupa dengan \texttt{plot3d}. Fungsi ini dapat memplot grafik
fungsi, matriks permukaan, maupun garis kontur opsional, lalu merender hasilnya secara
otomatis melalui Povray.

Selain menggunakan \texttt{pov3d()}, objek Povray kustom dapat dibuat secara eksplisit. Sesi
dimulai dengan perintah \texttt{povstart()}, diikuti penulisan kode Povray dengan \texttt{writeln()}, dan
diakhiri dengan \texttt{povend()} untuk memulai proses rendering.

Tekstur dan material objek dapat diatur menggunakan fungsi \texttt{povlook()}. Fungsi ini
mendukung pengaturan warna RGB, tingkat transparansi, pencahayaan ambien, serta efek
kilau permukaan (Phong shading).

Povray menggunakan sistem koordinat kidal (left-handed), sedangkan Euler menggunakan
sistem koordinat kanan (right-handed). EMT secara otomatis mengonversi koordinat sehingga
sumbu-z tetap mengarah ke atas secara vertikal.
Muat utilitas fungsional khusus pustaka Povray sebelum beroperasi.
\end{eulercomment}
\begin{eulerprompt}
>load povray;
\end{eulerprompt}
\begin{eulercomment}
Tentukan lokasi direktori instalasi Povray secara tepat. Sesuaikan jalur direktori berkas \texttt{pvengine.exe}
pada variabel \texttt{defaultpovray} jika diperlukan.
\end{eulercomment}
\begin{eulerprompt}
>defaultpovray="C:\Program Files\POV-Ray\v3.7\bin\pvengine.exe"
\end{eulerprompt}
\begin{euleroutput}
C:\Program Files\POV-Ray\v3.7\bin\pvengine.exe
\end{euleroutput}
\begin{eulercomment}
Contoh dasar perenderan kurva 3D dengan Povray. Proses ray tracing dijalankan di latar belakang.

Povray akan membuka jendela proses saat merender. Berikan izin akses jika dialog keamanan Windows
(UAC) muncul saat aplikasi dijalankan pertama kali.
\end{eulercomment}
\begin{eulerprompt}
>plot3d("x^2+y^2",zoom=2):
\end{eulerprompt}
\eulerimg{25}{images/Tugas 4-069.png}
\begin{eulerprompt}
>pov3d("x^2+y^2",zoom=3);
\end{eulerprompt}
\eulerimg{34}{images/Tugas 4-070.png}
\begin{eulercomment}
Properti tekstur, transparansi, serta garis tingkat kontur pada objek Povray dapat disesuaikan
menggunakan atribut pendukung.
\end{eulercomment}
\begin{eulerprompt}
>pov3d("x^2+y^3",axiscolor=red,angle=-45°,>anaglyph, ...
>look=povlook(cyan,0.2),level=-1:0.5:1,zoom=3.8);
\end{eulerprompt}
\eulerimg{37}{images/Tugas 4-071.png}
\begin{eulercomment}
Skala fungsi dapat disesuaikan secara manual agar bentuk plot sesuai dengan dimensi kubus
pembatas ketika fitur skala otomatis dimatikan (\texttt{<fscale}).

Contoh berikut memplot fungsi bidang kompleks dalam rentang koordinat \texttt{r=2}.
\end{eulercomment}
\begin{eulerprompt}
>pov3d("((x-1)^2+y^2)*((x+1)^2+y^2)/40",r=2, ...
>angle=-120°,level=1/40,dlevel=0.005,light=[-1,1,1],height=10°,n=50, ...
><fscale,zoom=3.8);
\end{eulerprompt}
\eulerimg{34}{images/Tugas 4-072.png}

\eulerheading{Memplot dengan Koordinat}
\begin{eulercomment}
Selain menggunakan ekspresi fungsi, objek Povray dapat dibentuk dari tiga matriks data
koordinat (x, y, z).

Contoh berikut menampilkan permukaan benda putar yang diputar mengelilingi sumbu-z:
\end{eulercomment}
\begin{eulerprompt}
>function f(x) := x^3-x+1; ...
>x=-1:0.01:1; t=linspace(0,2pi,50)'; ...
>Z=x; X=cos(t)*f(x); Y=sin(t)*f(x); ...
>pov3d(X,Y,Z,angle=40°,look=povlook(red,0.1),height=50°,axis=0,zoom=4,light=[10,5,15]);
\end{eulerprompt}
\eulerimg{34}{images/Tugas 4-073.png}
\begin{eulercomment}
Contoh berikut memperlihatkan gelombang terredam (damped wave) dalam koordinat polar.

Elemen tambahan dapat disisipkan ke dalam adegan Povray menggunakan parameter \texttt{add=}. Objek
tambahan tersebut secara otomatis disesuaikan ukurannya ke dalam skala kubus satuan.
\end{eulercomment}
\begin{eulerprompt}
>r=linspace(0,1,80); phi=linspace(0,2pi,80)'; ...
>x=r*cos(phi); y=r*sin(phi); z=exp(-5*r)*cos(8*pi*r)/3; ...
>pov3d(x,y,z,zoom=6,axis=0,height=30°,add=povsphere([0.5,0,0.25],0.15,povlook(red)), ...
>w=500,h=300);
\end{eulerprompt}
\eulerimg{24}{images/Tugas 4-074.png}
\begin{eulercomment}
Povray dapat menghasilkan permukaan yang sangat halus meskipun jumlah titik data grid
terbatas. Untuk mengoptimalkan tampilan shading dan bayangan, vektor normal permukaan
dapat disertakan dalam perhitungan.

Vektor normal dihitung dari turunan parsial permukaan fungsi:
\end{eulercomment}
\begin{eulerprompt}
>Z &= x^2*y^3
\end{eulerprompt}
\begin{euleroutput}
 2  3
x  y
\end{euleroutput}
\begin{eulercomment}
Perkalian silang (cross product) dari turunan terhadap x dan y memberikan vektor normal
tegak lurus permukaan.
\end{eulercomment}
\begin{eulerprompt}
>dx &= diff([x,y,Z],x); dy &= diff([x,y,Z],y);
\end{eulerprompt}
\begin{eulercomment}
Komponen vektor normal (NX, NY, NZ) dihitung secara simbolik:
\end{eulercomment}
\begin{eulerprompt}
>N &= crossproduct(dx,dy); NX &= N[1]; NY &= N[2]; NZ &= N[3]; N,
\end{eulerprompt}
\begin{euleroutput}
      3       2  2
[- 2 x y , - 3 x  y , 1]
\end{euleroutput}
\begin{eulercomment}
Meskipun hanya menggunakan grid 25x25 titik, permukaan terlihat halus karena
diterapkannya vektor normal (\texttt{xv}, \texttt{yv}, \texttt{zv}).
\end{eulercomment}
\begin{eulerprompt}
>x=-1:0.5:1; y=x';
>pov3d(x,y,Z(x,y),angle=10°, ...
>xv=NX(x,y),yv=NY(x,y),zv=NZ(x,y),<shadow);
\end{eulerprompt}
\eulerimg{34}{images/Tugas 4-075.png}
\begin{eulercomment}
Contoh simpul Trefoil (Trefoil Knot) berikut dirancang oleh A. Busser untuk EMT. Berkas
contoh ini dapat ditemukan di folder pustaka EMT.

Lokasi berkas: \texttt{Examples\textbackslash Trefoil Knot}.

Penggunaan vektor normal membuat permukaan kurva terlihat mulus tanpa memerlukan grid
yang terlalu padat. Pustaka Maxima digunakan untuk menghitung turunan parsial dan vektor
normal secara simbolik.
\end{eulercomment}
\begin{eulerprompt}
>X &= ((4+sin(3*y))+cos(x))*cos(2*y); ...
>Y &= ((4+sin(3*y))+cos(x))*sin(2*y); ...
>Z &= sin(x)+2*cos(3*y);
\end{eulerprompt}
\begin{eulercomment}
Turunan parsial terhadap x dan y dihitung secara simbolik:
\end{eulercomment}
\begin{eulerprompt}
>dx &= diff([X,Y,Z],x); dy &= diff([X,Y,Z],y);
\end{eulerprompt}
\begin{eulercomment}
Vektor normal hasil perkalian silang:
\end{eulercomment}
\begin{eulerprompt}
>dn &= crossproduct(dx,dy);
\end{eulerprompt}
\begin{eulercomment}
Evaluasi numerik dilakukan pada grid koordinat di EMT:
\end{eulercomment}
\begin{eulerprompt}
>x:=linspace(-%pi,%pi,40); y:=linspace(-%pi,%pi,100)';
\end{eulerprompt}
\begin{eulercomment}
Komponen vektor \texttt{dn} dievaluasi menggunakan sintaks \texttt{\&"expression"(parameters)}
untuk melewatkan ekspresi simbolik Maxima secara langsung ke \texttt{pov3d}.
\end{eulercomment}
\begin{eulerprompt}
>pov3d(X(x,y),Y(x,y),Z(x,y),>anaglyph,axis=0,zoom=5,w=450,h=350, ...
><shadow,look=povlook(blue), ...
>xv=&"dn[1]"(x,y), yv=&"dn[2]"(x,y), zv=&"dn[3]"(x,y));
\end{eulerprompt}
\eulerimg{28}{images/Tugas 4-076.png}
\begin{eulercomment}
Kerangka jaring (grid 3D) dapat dibuat dan di-render di Povray menggunakan perintah \texttt{povgrid()}.
\end{eulercomment}
\begin{eulerprompt}
>povstart(zoom=4); ...
>x=-1:0.5:1; r=1-(x+1)^2/6; ...
>t=(0°:30°:360°)'; y=r*cos(t); z=r*sin(t); ...
>writeln(povgrid(x,y,z,d=0.02,dballs=0.05)); ...
>povend();
\end{eulerprompt}
\eulerimg{34}{images/Tugas 4-077.png}
\begin{eulercomment}
Perintah \texttt{povgrid()} juga dapat digunakan untuk memplot kurva ruang 3D dalam bentuk tabung/kawat.
\end{eulercomment}
\begin{eulerprompt}
>povstart(center=[0,0,1],zoom=3.6); ...
>t=linspace(0,2,1000); r=exp(-t); ...
>x=cos(2*pi*10*t)*r; y=sin(2*pi*10*t)*r; z=t; ...
>writeln(povgrid(x,y,z,povlook(red))); ...
>writeAxis(0,2,axis=3); ...
>povend();
\end{eulerprompt}
\eulerimg{34}{images/Tugas 4-078.png}

\eulerheading{Objek Povray}
\begin{eulercomment}
Selain membuat grafik fungsi, Povray dapat membuat berbagai bentuk geometri murni dasar.
EMT menggenerasikan deskripsi teks objek tersebut ke berkas Povray.

Pembuatan adegan diawali dengan \texttt{povstart()}.
\end{eulercomment}
\begin{eulerprompt}
>povstart(zoom=4);
\end{eulerprompt}
\begin{eulercomment}
Contoh berikut mendefinisikan tiga silinder yang saling berpotongan tegak lurus.

Konstanta \texttt{povx}, \texttt{povy}, dan \texttt{povz} mewakili vektor satuan sumbu koordinat $[1,0,0]$, $[0,1,0]$,
dan $[0,0,1]$.
\end{eulercomment}
\begin{eulerprompt}
>c1=povcylinder(-povx,povx,1,povlook(red)); ...
>c2=povcylinder(-povy,povy,1,povlook(yellow)); ...
>c3=povcylinder(-povz,povz,1,povlook(blue)); ...
\end{eulerprompt}
\begin{eulercomment}
Kode Povray yang dihasilkan berupa teks berikut:
\end{eulercomment}
\begin{eulerprompt}
>c2
\end{eulerprompt}
\begin{euleroutput}
cylinder { <0,0,-1>, <0,0,1>, 1
    texture { pigment { color rgb <0.941176,0.941176,0.392157> }  } 
    finish { ambient 0.2 } 
}
\end{euleroutput}
\begin{eulercomment}
Tampilan tekstur dan warna tiap silinder diatur dengan fungsi \texttt{povlook()}.

Fungsi \texttt{povlook()} menghasilkan string kode tekstur Povray. Fungsi ini menerima masukan
warna RGB, transparansi, intensitas ambien, serta efek kilau permukaan.
\end{eulercomment}
\begin{eulerprompt}
>povlook(rgb(0.1,0.2,0.3),0.1,0.5)
\end{eulerprompt}
\begin{euleroutput}
texture { pigment { color rgbf <0.101961,0.2,0.301961,0.1> }  } 
finish { ambient 0.5 } 
\end{euleroutput}
\begin{eulercomment}
Operasi irisan (intersection) ketiga silinder digabungkan dengan perintah
\texttt{povintersection()}.
\end{eulercomment}
\begin{eulerprompt}
>writeln(povintersection([c1,c2,c3]));
\end{eulerprompt}
\begin{eulercomment}
Perpotongan ini menghasilkan bentuk geometri kompleks yang menarik.
\end{eulercomment}
\begin{eulerprompt}
>povend;
\end{eulerprompt}
\eulerimg{34}{images/Tugas 4-079.png}
\begin{eulercomment}
Objek fraktal 3D dapat dikonstruksi secara rekursif menggunakan fungsi bawaan.

Fungsi \texttt{povbox()} mengembalikan string definisi kubus Povray dengan koordinat dan tekstur
yang ditentukan.
\end{eulercomment}
\begin{eulerprompt}
>function onebox(x,y,z,d) := povbox([x,y,z],[x+d,y+d,z+d],povlook());
>function fractal (x,y,z,h,n) ...
\end{eulerprompt}
\begin{eulerudf}
if n==1 then writeln(onebox(x,y,z,h));
else
h=h/3;
fractal(x,y,z,h,n-1);
fractal(x+2*h,y,z,h,n-1);
fractal(x,y+2*h,z,h,n-1);
fractal(x,y,z+2*h,h,n-1);
fractal(x+2*h,y+2*h,z,h,n-1);
fractal(x+2*h,y,z+2*h,h,n-1);
fractal(x,y+2*h,z+2*h,h,n-1);
fractal(x+2*h,y+2*h,z+2*h,h,n-1);
fractal(x+h,y+h,z+h,h,n-1);
endif;
endfunction
\end{eulerudf}
\begin{eulerprompt}
>povstart(fade=10,<shadow);
>fractal(-1,-1,-1,2,4);
>povend();
\end{eulerprompt}
\eulerimg{34}{images/Tugas 4-080.png}
\begin{eulercomment}
Povray mendukung operasi geometri Boolean (Constructive Solid Geometry / CSG) seperti
penggabungan (union), irisan (intersection), dan pengurangan (difference) untuk memotong
atau menggabungkan objek geometri.
\end{eulercomment}
\begin{eulerprompt}
>povstart(light=[5,-5,5],fade=10);
\end{eulerprompt}
\begin{eulercomment}
Objek dapat didefinisikan sebagai variabel di Povray menggunakan perintah \texttt{povdefine()}.
Variabel ini ditulis langsung ke berkas Povray saat dipanggil.

Contoh berikut mendefinisikan kubus simetris bernilai -1 hingga 1:
\end{eulercomment}
\begin{eulerprompt}
>povdefine("mycube",povbox(-1,1));
\end{eulerprompt}
\begin{eulercomment}
Objek yang sudah didefinisikan dapat dipanggil kembali menggunakan \texttt{povobject()}.
\end{eulercomment}
\begin{eulerprompt}
>c1=povobject("mycube",povlook(red));
\end{eulerprompt}
\begin{eulercomment}
Kubus kedua diputar, digeser, dan diperbesar ukurannya:
\end{eulercomment}
\begin{eulerprompt}
>c2=povobject("mycube",povlook(yellow),translate=[1,1,1], ...
>rotate=xrotate(10°)+yrotate(10°), scale=1.2);
\end{eulerprompt}
\begin{eulercomment}
Fungsi \texttt{povdifference()} memotong kubus pertama dengan kubus kedua:
\end{eulercomment}
\begin{eulerprompt}
>writeln(povdifference(c1,c2));
\end{eulerprompt}
\begin{eulercomment}
Sumbu koordinat ditambahkan ke dalam adegan:
\end{eulercomment}
\begin{eulerprompt}
>writeAxis(-1.2,1.2,axis=1); ...
>writeAxis(-1.2,1.2,axis=2); ...
>writeAxis(-1.2,1.2,axis=4); ...
>povend();
\end{eulerprompt}
\eulerimg{34}{images/Tugas 4-081.png}

\eulerheading{Fungsi Implisit}
\begin{eulercomment}
Povray dapat merender permukaan fungsi implisit $f(x,y,z)=0$ secara langsung dengan
pencahayaan dan bayangan yang halus.

Penulisan fungsi implisit untuk Povray menggunakan sintaks bahasa C (misalnya \texttt{pow(x,2)}
untuk \texttt{x\textasciicircum 2}), bukan sintaks Euler atau Maxima.
\end{eulercomment}
\begin{eulerformula}
\[
((x^2+y^2-c^2)^2+(z^2-1)^2)*((y^2+z^2-c^2)^2+(x^2-1)^2)*((z^2+x^2-c^2)^2+(y^2-1)^2)=d
\]
\end{eulerformula}
\begin{eulerprompt}
>povstart(angle=70°,height=50°,zoom=4);
>writeln(povsurface("(pow(pow(x,2)+pow(y,2)-pow(0.1,2),2)+pow(pow(z,2)-1,2))*(pow(pow(y,2)+pow(z,2)-pow(0.1,2),2)+pow(pow(x,2)-1,2))*(pow(pow(z,2)+pow(x,2)-pow(0.1,2),2)+pow(pow(y,2)-1,2))-0.1",povlook(red)));
>povend();
\end{eulerprompt}
\eulerimg{34}{images/Tugas 4-083.png}
\begin{eulerprompt}
>povstart(angle=25°,height=10°);
>writeln(povsurface("pow(x,2)+pow(y,2)*pow(z,2)-1",povlook(blue),povbox(-2,2,"")));
>povend();
\end{eulerprompt}
\eulerimg{34}{images/Tugas 4-084.png}
\begin{eulerprompt}
>povstart(angle=70°,height=50°,zoom=4);
\end{eulerprompt}
\begin{eulercomment}
Fungsi \texttt{povsurface()} merender permukaan implisit menggunakan ekspresi bertipe C-style.
\end{eulercomment}
\begin{eulerprompt}
>writeln(povsurface("pow(x,2)*y-pow(y,3)-pow(z,2)",povlook(green))); ...
>writeAxes(); ...
>povend();
\end{eulerprompt}
\eulerimg{34}{images/Tugas 4-085.png}

\eulerheading{Objek Jaring (Mesh)}
\begin{eulercomment}
Objek jaring segi tiga (mesh object) dapat dibuat untuk menampilkan permukaan kompleks
beserta elemen dekoratif tambahan.

Contoh berikut memvisualisasikan maksimum fungsi $z = x \cdot y$ dengan batasan kondisi $x + y = 1$.
\end{eulercomment}
\begin{eulerprompt}
>povstart(angle=-10°,center=[0.5,0.5,0.5],zoom=7);
\end{eulerprompt}
\begin{eulercomment}
Data jaring yang besar disimpan ke berkas menggunakan perintah \texttt{\#declare}. Fungsi
\texttt{povtriangles()} menggenerasikan kode jaring segi tiga beserta vektor normalnya secara
otomatis.
\end{eulercomment}
\begin{eulerprompt}
>x=0:0.02:1; y=x'; z=x*y; vx=-y; vy=-x; vz=1;
>mesh=povtriangles(x,y,z,"",vx,vy,vz);
\end{eulerprompt}
\begin{eulercomment}
Dua objek cakram (disc) dibuat untuk memotong permukaan jaring menggunakan operasi CSG.
\end{eulercomment}
\begin{eulerprompt}
>cl=povdisc([0.5,0.5,0],[1,1,0],2); ...
>ll=povdisc([0,0,1/4],[0,0,1],2);
\end{eulerprompt}
\begin{eulercomment}
Bagian jaring di luar cakram dipotong menggunakan \texttt{povdifference()}.
\end{eulercomment}
\begin{eulerprompt}
>writeln(povdifference(mesh,povunion([cl,ll]),povlook(green)));
\end{eulerprompt}
\begin{eulercomment}
Irisan antara jaring dan cakram digambar dengan warna berbeda menggunakan
\texttt{povintersection()}.
\end{eulercomment}
\begin{eulerprompt}
>writeln(povintersection([mesh,cl],povlook(red))); ...
>writeln(povintersection([mesh,ll],povlook(gray)));
\end{eulerprompt}
\begin{eulercomment}
Titik maksimum fungsi ditandai dengan bola kecil menggunakan \texttt{povpoint()}.
\end{eulercomment}
\begin{eulerprompt}
>writeln(povpoint([1/2,1/2,1/4],povlook(gray),size=2*defaultpointsize));
\end{eulerprompt}
\begin{eulercomment}
Sumbu koordinat ditambahkan dengan \texttt{writeAxes()}, dan adegan diakhiri dengan \texttt{povend()}.
\end{eulercomment}
\begin{eulerprompt}
>writeAxes(0,1,0,1,0,1,d=0.015); ...
>povend();
\end{eulerprompt}
\eulerimg{34}{images/Tugas 4-086.png}

\eulerheading{Anaglyph dalam Povray}
\begin{eulercomment}
Efek anaglif 3D pada Povray dibuat dengan merender adegan dari dua posisi kamera yang
berbeda, lalu menggabungkan kedua gambar PNG menggunakan \texttt{loadanaglyph()}.

Gunakan kacamata filter merah/sian untuk melihat efek kedalaman 3D.

Mode anaglif pada \texttt{pov3d} dapat diaktifkan cukup dengan menambahkan sakelar \texttt{>anaglyph}.
\end{eulercomment}
\begin{eulerprompt}
>pov3d("-exp(-x^2-y^2)/2",r=2,height=45°,>anaglyph, ...
>center=[0,0,0.5],zoom=3.5);
\end{eulerprompt}
\eulerimg{37}{images/Tugas 4-087.png}
\begin{eulercomment}
Untuk adegan kustom, buat fungsi tersendiri yang mendefinisikan adegan. Fungsi tersebut
kemudian dipanggil oleh \texttt{povanaglyph()}.
\end{eulercomment}
\begin{eulerprompt}
>function myscene ...
\end{eulerprompt}
\begin{eulerudf}
s=povsphere(povc,1);
cl=povcylinder(-povz,povz,0.5);
clx=povobject(cl,rotate=xrotate(90°));
cly=povobject(cl,rotate=yrotate(90°));
c=povbox([-1,-1,0],1);
un=povunion([cl,clx,cly,c]);
obj=povdifference(s,un,povlook(red));
writeln(obj);
writeAxes();
endfunction
\end{eulerudf}
\begin{eulercomment}
Perintah \texttt{povanaglyph()} menjalankan pembuatan adegan sebanyak dua kali dari dua sudut
kamera berbeda untuk menghasilkan gambar anaglif.
\end{eulercomment}
\begin{eulerprompt}
>povanaglyph("myscene",zoom=4.5);
\end{eulerprompt}
\eulerimg{37}{images/Tugas 4-088.png}

\eulerheading{Mendefinisikan Objek Sendiri}
\begin{eulercomment}
Program Povray memiliki kemampuan untuk mencetak berbagai elemen spasial yang membentuk
geometri dasar. Modifikasi konstruksi objek juga dapat dilakukan tanpa batas.

Contoh berikut mendefinisikan bentuk torus (donat) di titik asal (0,0,0) menggunakan
fungsi kustom:
\end{eulercomment}
\begin{eulerprompt}
>function povdonat (r1,r2,look="") ...
\end{eulerprompt}
\begin{eulerudf}
return "torus {"+r1+","+r2+look+"}";
endfunction
\end{eulerudf}
\begin{eulercomment}
Torus pertama didefinisikan dengan jari-jari \texttt{r1=0.8} dan \texttt{r2=0.2}:
\end{eulercomment}
\begin{eulerprompt}
>t1=povdonat(0.8,0.2)
\end{eulerprompt}
\begin{euleroutput}
torus {0.8,0.2}
\end{euleroutput}
\begin{eulercomment}
Torus kedua diputar 90 derajat terhadap sumbu-x dan digeser ke sumbu-x sejauh 0.8:
\end{eulercomment}
\begin{eulerprompt}
>t2=povobject(t1,rotate=xrotate(90°),translate=[0.8,0,0])
\end{eulerprompt}
\begin{euleroutput}
object { torus {0.8,0.2}
    rotate 90 *x 
    translate <0.8,0,0>
}
\end{euleroutput}
\begin{eulercomment}
Kedua torus dimasukkan ke dalam adegan Povray dengan efek kilau \texttt{phong=1}:
\end{eulercomment}
\begin{eulerprompt}
>povstart(center=[0.4,0,0],angle=0°,zoom=3.8,aspect=1.5); ...
>writeln(povobject(t1,povlook(green,phong=1))); ...
>writeln(povobject(t2,povlook(green,phong=1))); ...
\end{eulerprompt}
\begin{eulerttcomment}
>povend();
\end{eulerttcomment}
\begin{eulercomment}
Jika terjadi kesalahan sintaks saat \texttt{povend()} mengeksekusi Povray, gunakan opsi \texttt{<exit}
agar jendela Povray tetap terbuka untuk memeriksa pesan kesalahan:
\end{eulercomment}
\begin{eulerttcomment}
>povend(<exit);
\end{eulerttcomment}
\begin{eulercomment}
Opsi ini mempertahankan jendela Povray terbuka jika terjadi galat.
\end{eulercomment}
\begin{eulerprompt}
>povend(h=320,w=480);
\end{eulerprompt}
\eulerimg{26}{images/Tugas 4-089.png}
\begin{eulercomment}
Contoh berikut memvisualisasikan daerah layak (feasible region) dan titik optimum masalah
pemrograman linear:
\end{eulercomment}
\begin{eulerformula}
\[
Ax \le b, \quad x \ge 0, \quad c.x \to \text{Max.}
\]
\end{eulerformula}
\begin{eulercomment}
Batas-batas ketaksamaan linear didefinisikan sebagai setengah-ruang (half-space).
\end{eulercomment}
\begin{eulerprompt}
>A=[10,8,4;5,6,8;6,3,2;9,5,6];
>b=[10,10,10,10]';
>c=[1,1,1];
\end{eulerprompt}
\begin{eulercomment}
Penyelesaian menggunakan algoritma simpleks mengembalikan titik optimum:
\end{eulercomment}
\begin{eulerprompt}
>x=simplex(A,b,c,>max,>check)'
\end{eulerprompt}
\begin{euleroutput}
[0,  1,  0.5]
\end{euleroutput}
\begin{eulercomment}
Setiap bidang pembatas didefinisikan dengan fungsi \texttt{oneplane()}:
\end{eulercomment}
\begin{eulerformula}
\[
a \cdot x \le b
\]
\end{eulerformula}
\begin{eulerprompt}
>function oneplane (a,b,look="") ...
\end{eulerprompt}
\begin{eulerudf}
return povplane(a,b,look)
endfunction
\end{eulerudf}
\begin{eulercomment}
Irisan dari semua bidang pembatas membentuk polihedron daerah layak:
\end{eulercomment}
\begin{eulerprompt}
>function adm (A, b, r, look="") ...
\end{eulerprompt}
\begin{eulerudf}
ol=[];
loop 1 to rows(A); ol=ol|oneplane(A[#],b[#]); end;
ol=ol|povbox([0,0,0],[r,r,r]);
return povintersection(ol,look);
endfunction
\end{eulerudf}
\begin{eulercomment}
Menampilkan adegan daerah layak beserta titik optimumnya di Povray:
\end{eulercomment}
\begin{eulerprompt}
>povstart(angle=120°,center=[0.5,0.5,0.5],zoom=3.5); ...
>writeln(adm(A,b,2,povlook(green,0.4))); ...
>writeAxes(0,1.3,0,1.6,0,1.5); ...
\end{eulerprompt}
\begin{eulercomment}
Bola merah ditambahkan untuk menandai lokasi titik optimum:
\end{eulercomment}
\begin{eulerprompt}
>writeln(povintersection([povsphere(x,0.5),povplane(c,c.x')], ...
>povlook(red,0.9)));
\end{eulerprompt}
\begin{eulercomment}
Vektor panah menunjukkan arah fungsi tujuan:
\end{eulercomment}
\begin{eulerprompt}
>writeln(povarrow(x,c*0.5,povlook(red)));
\end{eulerprompt}
\begin{eulercomment}
Teks 3D ditambahkan ke adegan menggunakan \texttt{povtext()}:
\end{eulercomment}
\begin{eulerprompt}
>writeln(povtext("Linear Problem",[0,0.2,1.3],size=0.05,rotate=5°)); ...
>povend();
\end{eulerprompt}
\eulerimg{34}{images/Tugas 4-092.png}

\eulerheading{Ringkasan Contoh Lainnya}
\begin{eulercomment}
Menggambar Grafik Fungsi Dua Variabel dalam Bentuk Ekspresi Langsung
\end{eulercomment}
\begin{eulerformula}
\[
f(x,y)=x^3-y^3
\]
\end{eulerformula}
\begin{eulerprompt}
>plot3d("x^3-y^3", -4, 4, -4, 4):
\end{eulerprompt}
\eulerimg{25}{images/Tugas 4-094.png}
\begin{eulercomment}
Menggambar Grafik Fungsi Dua Variabel yang Rumusnya Disimpan dalam Variabel Ekspresi
\end{eulercomment}
\begin{eulerformula}
\[
f(x,y)=x^2sin(y)
\]
\end{eulerformula}
\begin{eulerprompt}
>expr="x^2*sin(y)";
>plot3d(expr,-10,10,0,4*pi):
\end{eulerprompt}
\eulerimg{25}{images/Tugas 4-096.png}
\begin{eulercomment}
Menggambar Grafik Fungsi Dua Variabel yang Fungsinya Didefinisikan sebagai Fungsi Numerik
\end{eulercomment}
\begin{eulerformula}
\[
f(x,y)=x^2+y^2
\]
\end{eulerformula}
\begin{eulerprompt}
>function f(x,y):=x^2+y^2
>plot3d("f",-3,3,-3,3):
\end{eulerprompt}
\eulerimg{25}{images/Tugas 4-098.png}
\begin{eulercomment}
Menggambar Grafik Fungsi Dua Variabel yang Fungsinya Didefinisikan sebagai Fungsi Simbolik
\end{eulercomment}
\begin{eulerformula}
\[
f(x,y)=cos(x)sin(y)
\]
\end{eulerformula}
\begin{eulerprompt}
>function f(x,y)&=cos(x)*sin(y)
\end{eulerprompt}
\begin{euleroutput}
cos(x) sin(y)
\end{euleroutput}
\begin{eulerprompt}
>plot3d("f",-pi,pi,-pi,pi):
\end{eulerprompt}
\eulerimg{25}{images/Tugas 4-100.png}
\begin{eulercomment}
Menggambar Data x, y, z pada ruang Tiga Dimensi (3D)
\end{eulercomment}
\begin{eulerprompt}
>x=linspace(-3,3,40)';
>y=linspace(-3,3,40);
>z=exp(-(x^2+y^2));
>...plot3d(x,y,z):
\end{eulerprompt}
\begin{eulercomment}
Membuat Gambar Grafik Tiga Dimensi (3D) yang Bersifat Interaktif dan Animasi Grafik 3D
\end{eulercomment}
\begin{eulerformula}
\[
z(x,y,t)=sin(x+t)cos(y)
\]
\end{eulerformula}
\begin{eulerprompt}
>clg; for t=0 to 2*pi step 0.2; plot3d("sin(x+t)*cos(y)", -pi, pi, -pi, pi, r=2); wait(0.05); end;
\end{eulerprompt}
\begin{eulercomment}
Menggambar Fungsi Parametrik Tiga Dimensi (3D)
\end{eulercomment}
\begin{eulerformula}
\[
x(x,y)=cos(x)cos(y)
\]
\end{eulerformula}
\begin{eulerformula}
\[
y(x,y)=sin(x)cos(y)
\]
\end{eulerformula}
\begin{eulerformula}
\[
z(x,y)=sin(y)
\]
\end{eulerformula}
\begin{eulercomment}
dengan domain \texttt{x} pada $[0,2\pi]$ dan \texttt{y} pada $[-2\pi,2\pi]$
\end{eulercomment}
\begin{eulerprompt}
>plot3d("cos(x)*cos(y)", "sin(x)*cos(y)", "sin(y)", 0, 2*pi, -pi/2, pi/2):
\end{eulerprompt}
\eulerimg{25}{images/Tugas 4-105.png}
\begin{eulercomment}
Menggambar Fungsi Implisit Tiga Dimensi (3D)
\end{eulercomment}
\begin{eulerformula}
\[
x^2+y^2+z^2=4
\]
\end{eulerformula}
\begin{eulerprompt}
>function f(x,y,z)&= x^2+y^2+z^2-4;
>plot3d("f", >implicit, r=3):
\end{eulerprompt}
\eulerimg{25}{images/Tugas 4-107.png}
\begin{eulercomment}
Menggambar Titik pada Ruang Tiga Dimensi (3D)
\end{eulercomment}
\begin{eulerprompt}
>t=linspace(0,4*pi,80);
>x=cos(t); y=sin(t); z=t;
>plot3d(x,y,z, points=true, >wire):
\end{eulerprompt}
\eulerimg{25}{images/Tugas 4-108.png}
\begin{eulercomment}
Mengatur Tampilan, Warna, dan Sudut Pandang Gambar Permukaan Tiga Dimensi (3D)
\end{eulercomment}
\begin{eulerformula}
\[
f(x,y)=sin(x)cos(y)
\]
\end{eulerformula}
\begin{eulerprompt}
>plot3d("sin(x)*cos(y)", -pi, pi, -pi, pi, angle=45, height=35, >hue):
\end{eulerprompt}
\eulerimg{25}{images/Tugas 4-110.png}
\begin{eulercomment}
Menampilkan Kontur dan Bidang Kontur Permukaan Tiga Dimensi (3D)
\end{eulercomment}
\begin{eulerformula}
\[
f(x,y)=x^2+y^2
\]
\end{eulerformula}
\begin{eulerprompt}
>plot3d("x^2 + y^2", -2, 2, -2, 2, >contour, >cp):
\end{eulerprompt}
\eulerimg{25}{images/Tugas 4-112.png}
\begin{eulercomment}
Menggambar Grafik Tiga Dimensi dalam Modus Anaglif
\end{eulercomment}
\begin{eulerformula}
\[
f(x,y)=sin(x)cos(y)
\]
\end{eulerformula}
\begin{eulerprompt}
>plot3d("sin(x)*cos(y)", -pi, pi, -pi, pi, >anaglyph):
\end{eulerprompt}
\eulerimg{25}{images/Tugas 4-114.png}
\begin{eulercomment}
Menggambar Diagram Batang Tiga Dimensi
\end{eulercomment}
\begin{eulerprompt}
>M=[1, 3, 2, 5; 4, 2, 6, 3; 2, 5, 1, 4];
>plot3d(M, >bar):
\end{eulerprompt}
\eulerimg{25}{images/Tugas 4-115.png}
\begin{eulercomment}
Menggambar Permukaan Benda Putar
\end{eulercomment}
\begin{eulerprompt}
>clg; plot2d(0, 0, a=0, b=10, c=0, d=10, grid=1, title="Klik 5 titik"); ...
>x = []; y = []; ...
>for i=1 to 5; p=mouse(); x=[x,p[1]]; y=[y,p[2]]; plot2d(x, y, a=0, b=10, c=0, d=10, >add, style="-", color=blue); plot2d(x, y, a=0, b=10, c=0, d=10, >points, color=red); wait(0); end; ...
>plot2d(x, y, >add, style="o-", color=blue); wait(2); ...
>v = linspace(0, 2*pi, 60); X = x'; Y = y'*cos(v); Z = y'*sin(v); plot3d(X, Y, Z, frame=3, angle=40°, height=25°, >user):
\end{eulerprompt}
\eulerimg{25}{images/Tugas 4-116.png}
\begin{eulercomment}
Menggambar Grafik 3D dengan Povray di EMT
\end{eulercomment}
\begin{eulerprompt}
>povstart(angle=40°, height=25°, zoom=3); ...
>pov3d(X, Y, Z); ...
>povend();
\end{eulerprompt}
\eulerimg{34}{images/Tugas 4-117.png}
\eulerimg{34}{images/Tugas 4-118.png}

\eulerheading{Referensi Contoh Tambahan}
\begin{eulercomment}
Berkas panduan dan contoh pemodelan Povray tambahan dapat ditemukan pada direktori instalasi EMT
berikut:
\begin{itemize}
    \item \texttt{Examples/Dandelin Spheres}
    \item \texttt{Examples/Donat Math}
    \item \texttt{Examples/Trefoil Knot}
    \item \texttt{Examples/Optimization by Affine Scaling}
\end{itemize}
\end{eulercomment}

\end{eulernotebook}
\end{document}